\section{\texorpdfstring{\(p\) 值不是原假设为真的概率}{p 值不是原假设为真的概率}}
\label{sec:11-Mathematical-Statistics/05-Hypothesis-Testing/02}

你刚跑完一轮实验，\(p = 0.03\)。很自然地，你对合作者说："新方案有效的概率是 97\%。"

停一下。\(p\) 值的定义里，找不到"原假设为真的概率"这个量。它算的是另一个方向的东西。

\subsection{一个数字化困境：\(p\) 值到底在算什么}

设想一位研究者每收集 10 个样本就检查一次 \(p\) 值，一旦小于 \(0.05\) 就停止实验。在全部原假设都为真的重复操作中，最终"至少一次显著"的概率不是 5\%，而是远高于此——可能飙到 20\% 甚至更高。两种操作的 \(p\) 值计算公式完全相同，但错误率完全不同。问题不在计算，而在分析时点没有预先固定。

先把操作分歧放一边，回到定义本身。选定检验统计量 \(T(X)\)，且约定越大的 \(T\) 越不利于 \(H_0\)。观察到 \(t_{\mathrm{obs}}\) 后，简单原假设下的 \(p\) 值定义为

\[
p = P_{H_0}
(T(X)\ge t_{\mathrm{obs}}).
\]

把公式翻译成操作语言：假设原假设成立，按完全相同的实验设计重新采样，得到至少和当前观察同样极端的统计量的概率。

这个定义有三个关键点，缺一个都会跑偏。第一，"同样极端"的方向必须预先规定——先定好尾部方向，再算概率，不能看到数据后临时挑一个让自己更显著的方向。第二，概率完全计算在 \(H_0\) 的地盘上，\(H_0\) 是前提，不是结论。第三，\(p\) 值是一个关于数据的概率陈述，不是关于假设的概率陈述。双侧检验中，不能总是机械地把单侧 \(p\) 值乘二；非对称离散分布里，"同样或更极端"的排序需要更细致的约定。

\subsection{碰到真正的墙：条件方向反了}

从这里开始，真正的分歧出现了。很多人读完定义后仍然下意识地认为"\(p = 0.03\) 意味着原假设为真的概率是 3\%"。这个直觉错了，但不是粗心的问题，而是条件概率的方向问题。

\(p\) 值计算的是

\[
P(\text{数据}\ge\text{当前极端程度}\given H_0\text{为真}),
\]

而研究者想知道的是

\[
P(H_0\text{为真}\given\text{观测到的数据}).
\]

两者之间隔着 Bayes 公式：

\[
P(H_0\given\text{数据})
=
\frac{P(\text{数据}\given H_0)\,P(H_0)}
{P(\text{数据})}.
\]

\(p\) 值只提供了分子中的第一项 \(P(\text{数据}\given H_0)\)，还缺先验概率 \(P(H_0)\) 和边缘似然 \(P(\text{数据})\)。没有这两个量，无法从数据跨回假设。\(p\) 值定义中从未出现过先验分布，它根本没有装配这个功能。

用医学检验打个比方：假定你没生病，查出阳性指标的概率可能是 5\%（这就是 \(p\) 值的地盘）；但你拿到阳性报告后实际生病的概率是多少？那取决于疾病的基准患病率和检验的特异度。\(p = 0.03\) 和"生病的概率是 3\%"之间的差距，和这个例子一模一样。

\subsection{复合原假设要控制最不利参数}

上面讨论的是简单原假设——\(H_0\) 只包含一个分布。实际中更常见的是复合原假设

\[
H_0:\theta\in\Theta_0,
\]

包含多个可能的参数值。此时一个合法构造是取上确界：

\[
p = \sup_{\theta\in\Theta_0}
P_\theta
(T(X)\ge t_{\mathrm{obs}}).
\]

这保证了无论 \(\theta\) 在 \(\Theta_0\) 中取什么值，第一类错误都不会超过宣称的水平。如果你嫌上确界太保守，挑一个"方便"的参数值代替最不利值，得到的 \(p\) 值可能会反保守——它宣称的显著性比真实情况更夸张。

好在对于连续简单原假设且统计量分布连续的情形，\(p\) 值在 \(H_0\) 下服从 \(\mathrm{Uniform}(0,1)\)。这个性质来自概率积分变换：把连续随机变量代入自己的分布函数，结果就是均匀分布。这意味着如果你一辈子只做真原假设的检验，你的 \(p\) 值会在 \([0,1]\) 上均匀铺开，不会往 0 附近集中。离散检验的 \(p\) 值通常更保守，分布不再连续均匀，不会覆盖 \([0,1]\) 上的每一个数值。

\subsection{一个具体的 \(z\) 检验}

取双侧标准正态检验，观察到

\[
|z_{\mathrm{obs}}|=2.
\]

标准正态的对称性让双侧 \(p\) 值恰好等于单侧尾概率的两倍：

\[
p = 2P(Z\ge2)
\approx0.0455.
\]

在预先固定的 \(\alpha = 0.05\) 检验中会拒绝 \(H_0\)，在 \(\alpha = 0.01\) 中不会。\(p\) 值可以理解为"使当前结果刚好进入拒绝域的最小常用显著性水平"，这个解读在技术上没问题，但它容易让人误以为 \(p\) 值是一个连续的"可信度刻度"——一个研究结论到底有多可信，看 \(p\) 值大小就够了。

数字的陷阱就在这里。\(p = 0.049\) 和 \(p = 0.051\) 在证据意义上几乎相同，却会被 \(\alpha = 0.05\) 的二元标签切成"显著"和"不显著"。报告精确 \(p\) 值、效应估计值和置信区间，比只写一个"显著/不显著"的标签传递更多信息。

\subsection{\(p\) 值不回答哪些问题}

\(p = 0.03\) 不意味着：

\begin{itemize}
\tightlist
\item
  \(H_0\) 为真的概率是 3\%；
\item
  结果由随机机会造成的概率是 3\%；
\item
  \(H_1\) 为真的概率是 97\%；
\item
  复现实验得到显著结果的概率是 97\%；
\item
  效应很大或具有实际意义；
\item
  数据、模型和研究设计没有偏差。
\end{itemize}

每一条背后都有一个数学原因。要给假设本身分配概率，需要先验与 Bayesian 模型——\(p\) 值定义里从来没有先验。要讨论复现概率，需要指定真实效应的大小、未来的实验设计和分析流程——\(p\) 值只用到了当前这批数据。要判断效应是否重要，需要效应尺度和决策损失函数——\(p\) 值只关心"是否非零"，不关心"偏离多大"。

先验与似然比如何合成后验概率，下一节\hyperref[sec:11-Mathematical-Statistics/05-Hypothesis-Testing/03]{``似然比不是后验赔率''}会具体展开。那里的核心信息是：即使你用了似然比，没有先验也无法得到后验概率——似然比只是 Bayes 公式中的一项，和 \(p\) 值处于同一个方向，缺先验就无法翻身。

\subsection{可选停止会破坏尾概率校准}

回到开头那个困境。若研究者每收集一些样本就计算一次 \(p\) 值，一旦小于 \(0.05\) 就停止，最终"至少一次显著"的概率会远超 \(0.05\)。这里的问题不是某一次计算错了，而是整个分析流程违背了固定样本 \(p\) 值的预设——它假设分析时点在实验开始前就定好了，中间不能因为看到数据而改变决定。

合法的序贯检验（sequential test）、alpha spending 函数、always-valid \(p\) 值或 e-value 可以在中途查看数据，但必须从设计阶段就采用相应的理论框架。不能实验做完后，把固定样本检验的公式套在一系列中途查看上，然后宣称结果显著。

类似的陷阱也出现在分析流程的其他环节：先尝试多种排除规则、多种协变量、多种终点和多种模型，最后只报告那个最小的 \(p\) 值。这种操作相当于隐形的多重比较——表面上只做了一次检验，实际上大脑已经看过几十个结果并挑出了最有利的一个。其计量与校正见\hyperref[sec:11-Mathematical-Statistics/05-Hypothesis-Testing/04]{``多重比较改变错误的计量单位''}一节。

\subsection{\(p\) 值与置信区间的联系}

双侧 level-\(\alpha\) 检验与 \(1-\alpha\) 置信区间通常互为反演：

\[
H_0:\theta=\theta_0
\text{ 被拒绝}
\quad\Longleftrightarrow\quad
\theta_0\notin C_{1-\alpha}(X).
\]

这个对应关系意味着区间和检验在拒绝/不拒绝层面上是等价的。但区间能额外显示效应的大小和精度——你不仅知道"是否为零"，还能看到估计值落在什么范围内。因此报告区间比只报告一个阈值化的 \(p\) 值更有信息量。

不过要小心：若区间方法和检验方法不同——例如一个用 Wald 近似、另一个用 likelihood ratio——数值上可能不完全对应。反演等价性依赖于两者基于同一种构造逻辑。换了构造方法，区间端点可能不再恰好等于检验的拒绝边界。

回到那句话：\(p = 0.03\) 不意味着新方案有 97\% 的概率有效。它只意味着：如果新方案完全无效，你看到当前这么极端数据的概率是 3\%。这个信息有用，但它的含义到此为止，没有多跨一步的权利。
